% !TEX root = ../McCoy-Dissertation.tex
\chapter{Conclusions and Outlook}\label{ch:conclusions}
%%%%%%%%%%%%%%%%%%%%%%%%%%%%%%%%%%%%%%%%%--------------------------------------------------
This dissertation contributes to the field of high-energy astrophysics through technological development of specialized reflection gratings that push the state of the art for x-ray spectroscopy. 
A main scientific goal for the currently-planned \emph{Lynx X-ray Observatory} \cite{Gaskin19} is to diagnose highly-ionized portions of extended galactic halos and intergalactic medium by measuring the equivalent widths of weak absorption lines produced from hydrogen-like and helium-like species of oxygen, and other $\mathcal{Z} \geq 6$ elements [\emph{cf.\@} \cref{tab:astro_element_nuclei_mass}], along the line-of-sight of active galactic nuclei [\emph{cf.\@} \cref{sec:astro_plasmas}].  
The detection of these weak spectral lines, which fall in the soft x-ray bandpass [\emph{cf.\@} \cref{tab:astro_element_lyman_approx,tab:astro_element_he_res}], relies not only on having a large instrument collecting area for spectral sensitivity, $A_{\text{col}}$, but also on achieving high spectral resolving power, $\mathscr{R}$, with a figure of merit depending on $\sqrt{A_{\text{col}} \mathscr{R}}$ [\emph{cf.\@} \cref{eq:FOM_weak}]. 
As discussed in \cref{sec:grating_tech_dev,sec:ch1_conclusions}, implementing reflection gratings for the \emph{X-ray Grating Spectrometer (XGS)} planned for \emph{Lynx} \cite{McEntaffer19} centers on
\begin{enumerate}[noitemsep]
	\item the production of a master grating with a sawtooth surface-relief profile that enables 
	\begin{enumerate}[noitemsep]
    	\item total absolute diffraction efficiency in the soft x-ray exceeding \SI{40}{\percent}, and 
        \item $\mathscr{R} \gtrapprox 5000$ in a Wolter-I telescope
    \end{enumerate}
	and additionally,  
	\item the mass manufacture of grating replicas to populate modular grating arrays, which involves 
	\begin{enumerate}[noitemsep]
    	\item imprinting grating surface-relief molds, and 
        \item coating each replica for soft x-ray reflectivity.  
    \end{enumerate}
\end{enumerate}
With these requirements motivated in \cref{ch:introduction}, \cref{ch:master_fab,ch:grating_replication} focus on applying the processes of \emph{thermally-activated selective topography equilibration (TASTE)} \cite{Schleunitz14} and \emph{substrate-conformal imprint lithography (SCIL)} \cite{Verschuuren17} to x-ray reflection grating technology, with an emphasis on characterizing diffraction efficiency using the beamline methodology described in \cref{ch:diff_eff}. 
The findings from these studies on TASTE and SCIL are summarized in \cref{sec:TASTE_ch5,sec:SCIL_ch5}, respectively, with future work described therein. 
Finally, outlook for potential future studies is discussed in \cref{sec:summary_ch5}.
%outlook for future work provided therein. 
%Finally, a summary and a discussion on future work is provided in \cref{sec:summary_ch5}.

\section{TASTE for Master Grating Fabrication}\label{sec:TASTE_ch5}
%%%%%%%%%%%%%%%%%%%%%%%%%%%%%%%%%%%%%%%%%--------------------------------------------------
Coupling \emph{grayscale electron-beam lithography (GEBL)} with selective \emph{thermal reflow} \cite{Schleunitz10,Schleunitz11a,Schleunitz14,Kirchner19}, TASTE offers an avenue for patterning sawtooth-like topographies in polymeric resist over a custom groove layout. 
%Utilizing this approach for 
This approach to grating fabrication eliminates dependence on substrate crystal structure, thereby providing a way for a radially-ruled, blazed grating layout to be generated with the precision of EBL, which, in principle, enables high $\mathscr{R}$ by reducing spectral aberrations from an imperfect radial profile [\emph{cf.\@} \cref{sec:crystal_etching}]. 
Through carrying out process development for TASTE using the \textsc{EBPG5200} tool at the Penn State Nanofabrication Laboratory \cite[\emph{cf.\@} \cref{fig:ebeam_tool_pic}]{PSU_MRI_EBL,raith_ebpg}, it was found that sub-\si{\um}, sawtooth-like topographies could be generated in \SI{130}{\nm}-thick poly(methyl methacrylate) (PMMA) [\emph{cf.\@} \cref{fig:PMMA}] by patterning repeating staircase patterns by GEBL and then performing an appropriate thermal annealing step. 
The results presented in \cref{sec:TASTE} show a path forward for fabricating x-ray reflection gratings using TASTE. 
Patterns with a groove spacing of $d = \SI{840}{\nm}$ yield a blaze angle $\delta \approx 10^{\circ}$ and are more suitable for in-plane geometries while patterns with $d = \SI{400}{\nm}$ provide a base for off-plane gratings with $\delta \approx 25^{\circ}$ \cite{McCoy18}. 

Starting with a $d = \SI{400}{\nm}$ test pattern from \cref{sec:TASTE}, a grating prototype was fabricated by coating TASTE-processed PMMA on a silicon wafer with a thin layer of gold for reflectivity by \emph{electron-beam physical vapor deposition (EBPVD)}, using titanium as an adhesion layer [\emph{cf.\@} \cref{sec:TASTE_prototype}]. 
This functional reflection grating was then tested for diffraction efficiency in an extreme off-plane mount over the photon energy range \SIrange{80}{800}{\electronvolt} at beamline 6.3.2 of the Advanced Light Source (ALS) at Lawrence-Berkeley National Laboratory \cite{ALS_632,Underwood96,Gullikson01,LBNL_web}. 
%As described in \cref{sec:taste_eff_results}, 
The results show absolute, peak-order efficiency ranging from \SIrange{75}{25}{\percent} as photon energy, $\mathcal{E}_{\gamma}$ increases while total diffraction efficiency, $\mathscr{E}_{\text{tot}}$, is maintained at $\lessapprox \SI{60}{\percent}$ across the tested bandpass, which demonstrates that x-ray reflection gratings fabricated by TASTE are capable of meeting \emph{Lynx} requirements in terms of spectral sensitivity \cite[\emph{cf.\@} \cref{sec:taste_eff_results}]{McCoy20}. 
\begin{table}[]
\centering
\caption[Benchmark table of blazed, off-plane x-ray reflection gratings in terms of total diffraction efficiency near the L-shell resonance transition in \ion{O}{vii}]{Benchmark table of blazed, off-plane x-ray reflection gratings in terms of \emph{total absolute diffraction efficiency} ($\mathscr{E}_{\text{tot}} \equiv \sum_n \mathscr{E}_n$ for all propagating orders with $n \neq 0$) near the L-shell resonance transition in \ion{O}{vii} at $\mathcal{E}_{\gamma} \approx \SI{574}{\electronvolt}$ [\emph{cf.\@} \cref{tab:astro_element_he_res}]. Refs~\cite{Werner77,Cash82,McEntaffer04} use data gathered at discrete energies using electron-impact sources; in these cases $\mathscr{E}_{\text{tot}}$ is referenced at the $\mathcal{E}_{\gamma} \approx \SI{525}{\electronvolt}$ K-shell fluorescence line in neutral oxygen. Refs~\cite{Tutt16,DeRoo16thesis,Miles18,Miles19,McCoy20,McCurdy20} use data gathered at synchrotron facilities; in these cases $\mathscr{E}_{\text{tot}}$ is referenced using measurements closest to $\mathcal{E}_{\gamma} \approx \SI{574}{\electronvolt}$ (typically within \SI{25}{\electronvolt}). For comparison, critical transmission gratings yield $\mathscr{E}_{\text{tot}} \approx \SI{30}{\percent}$ near  $\mathcal{E}_{\gamma} \approx \SI{574}{\electronvolt}$ \cite{Heilmann11,Heilmann16}. Items are listed in order of publication year (1977-2020).}\label{tab:eff_benchmark}
\begin{tabular}{cccc}
\cline{1-4}
fabrication           & grating grooves & geometry & $\mathscr{E}_{\text{tot}}$ \\ \hline
mechanical ruling \cite{Werner77} & $\delta \approx 5^{\circ}$, $d \approx \SI{278}{\nm}$ (\ce{Au}) & $\alpha \approx 0$, $\gamma \lessapprox 2^{\circ}$   & $\sim \SI{20}{\percent}$            \\
mechanical ruling \cite{Cash82}   & $\delta \approx 21^{\circ}$, $d \approx \SI{167}{\nm}$ (\ce{Au})  & $\alpha \approx \delta$, $\gamma \lessapprox 2^{\circ}$  & $\sim \SI{25}{\percent}$           \\
holography/ion etch \cite{McEntaffer04} & $\delta \approx 9^{\circ}$, $d \approx \SI{236}{\nm}$ (\ce{Pt})  & $\alpha \approx \delta$, $\gamma \approx 2.7^{\circ}$      & \SI{27}{\percent}            \\
holography/ion etch \cite{Tutt16} & $\delta \approx 9^{\circ}$, $d \approx \SI{236}{\nm}$ (\ce{Pt})  & $\alpha \approx 0$, $\gamma \approx 2^{\circ}$ & $\sim \SI{25}{\percent}$  \\
holography/ion etch \cite{Tutt16} & $\delta \approx 16^{\circ}$, $d \approx \SI{170}{\nm}$ (\ce{Pt})  & $\alpha \approx 0$, $\gamma \approx 2^{\circ}$ & $\sim \SI{30}{\percent}$  \\
\ce{KOH} etch (via EBL) \cite{DeRoo16thesis}  & $\delta \lessapprox 55^{\circ}$, $d = \SI{160}{\nm}$ (\ce{Au})  & $\alpha \approx  \delta$, $\gamma \approx 2.6^{\circ}$  & $\gtrapprox \SI{35}{\percent}$  \\
\ce{KOH} etch/UV-NIL \cite{Miles18}  & $\delta \lessapprox 30^{\circ}$, $d \lessapprox \SI{160}{\nm}$ (\ce{Au})  & $\alpha \approx 25^{\circ}$, $\gamma \approx 1.7^{\circ}$         & $\lessapprox \SI{60}{\percent}$  \\
\ce{KOH} etch/SCIL \cite{Miles19}  & $\delta \lessapprox 30^{\circ}$, $d \lessapprox \SI{174}{\nm}$ (\ce{Ni})  & $\alpha \approx 23^{\circ}$, $\gamma \approx 2.5^{\circ}$         & $\sim \SI{55}{\percent}$  \\
TASTE \cite[\emph{cf.\@} \cref{ch:master_fab}]{McCoy20}  & $\delta \approx 27^{\circ}$, $d = \SI{400}{\nm}$ (\ce{Au})  &  $\alpha \approx 23^{\circ}$, $\gamma \approx 1.7^{\circ}$        & $\lessapprox \SI{60}{\percent}$            \\
TASTE \cite{McCurdy20}  & $\delta \approx 35^{\circ}$, $d = \SI{160}{\nm}$ (\ce{Au}) & $\alpha \approx \delta$, $\gamma \approx 1.8^{\circ}$         & \SI{43}{\percent}  \\ \hline
\end{tabular}
\end{table}
This measurement for $\mathscr{E}_{\text{tot}}$ is compared to analogous results from other blazed x-ray reflection gratings in \cref{tab:eff_benchmark}, where it can be seen that \ce{KOH}-etched gratings and TASTE-fabricated gratings offer substantial improvement over gratings manufactured by mechanical ruling engine or by holographic recording coupled with directional ion etching [\emph{cf.\@} \cref{sec:grating_tech_dev}]. 

%Because this latter condition relies on the realization of a radial groove profile that matches the focal length of the telescope, beamline testing for $\mathscr{R}$ is required to make this determination. 
%Achieving a high-fidelity radial profile in part requires careful control of field stitching with the \textsc{EBPG5200} especially along the grating-dispersion direction so that irregularities in $d$, which can lead to a degradation in $\mathscr{R}$ as alluded to in \cref{sec:binary_ebeam}, can be avoided. 

While the prototype diffraction-efficiency results presented in \cref{sec:taste_eff_results} have been shown to be characteristic of a blazed grating with $\delta \approx 27^{\circ}$ through \textsc{PCGrate-SX} modeling \cite[\emph{cf.\@} \cref{sec:integral_method}]{PCGrate_web,Goray10}, there exists room for improvement in this fabrication process. 
In particular, the fidelity of the sawtooth-like groove facets hinges on the selectivity of the thermal reflow step, which in turn depends on the distribution of average molecular weight in the resist, $M_w$, imparted via high-energy electron exposure in GEBL [\emph{cf.\@} \cref{fig:TASTE_diagram}]. 
This material contrast in $M_w$ establishes a lateral gradient in the polymer \emph{glass transition temperature}, $T_g$, such that the viscosity of molten resist at a given temperature decreases with $M_w$. 
Selective thermal reflow is then induced by heating the resist by hotplate to a carefully chosen temperature, $T_{\text{reflow}}$, that is lower than $T_g$ for unexposed portions of the resist (\emph{i.e.}, top steps in a staircase pattern), but higher than the range of $T_g$ associated with electron-exposed resist (\emph{i.e.}, the remaining, lower steps in a staircase pattern) \cite{Schleunitz14,Kirchner19}. 
Because resist can be dosed inadvertently from electron scattering, especially when the widths of staircase steps are comparable to the size of the focused electron beam (typically $\lessapprox \SI{40}{\nm}$ in the \textsc{EBPG5200}), the size of this process window in practice depends on critical dimensions of the GEBL layout. 
Thermal reflow selectivity in the $d = \SI{400}{\nm}$ grating prototype described in \cref{sec:TASTE_prototype} appears to be diminished relative to \si{\um}-scale TASTE patterns previously reported in the literature \cite{Schleunitz10,Schleunitz14,Kirchner16,Pfirrmann16}, which is evidenced by the slight bulging effect seen in \cref{fig:B_time_scan,fig:TASTE_AFMs}.  
Moreover, McCurdy, et~al.~\cite{McCurdy20} show that this effect is exacerbated at $d = \SI{160}{\nm}$, where a decreased selectivity for thermal reflow causes significant rounding in the top staircase step as well as facet-surface irregularities that ultimately cause lower overall diffraction efficiency compared to what is reported in \cref{sec:taste_eff_results} [\emph{cf.\@} \cref{tab:eff_benchmark}]. %arising from the GEBL layout 

Other limitations of TASTE for grating fabrication stem from its dependence on resist thickness for a given GEBL recipe. 
Beyond the need for a uniformly-thick layer of resist, which can be difficult to achieve consistently through spin-coating, the groove depth (\emph{i.e.}, the height of top steps in a staircase pattern) is determined essentially by the resist thickness. 
This then places a restriction on $\delta$ for a fixed value of $d$, where $\delta$ decreases with increasing $d$; for example, the results from \cref{sec:TASTE} show that groove spacings of $d = \SI{400}{\nm}$ and $d = \SI{840}{\nm}$ at a resist thickness of \SI{130}{\nm} yield $\delta \approx 25^{\circ}$ and $\delta \approx 10^{\circ}$, respectively. 
Changing the resist thickness for GEBL requires the establishment of a new resist-contrast curve [\emph{cf.\@} \cref{sec:resist_contrast}]. 
For any given resist thickness, however, the blaze angle is expected to be limited practically to $\delta \lessapprox 45^{\circ}$ due to the difficulties involved with patterning high-aspect-ratio structures in GEBL. 
On the other hand, a TASTE pattern can be transferred into silicon through a dry etch, and depending on the etch selectivity of the resist relative to the substrate, $\delta$ can, in principle, be steepened relative to what is patterned directly in resist. 
A desired etch selectivity may be achieved by using a different resist (\emph{e.g.}, ZEP520A, mr-PosEBR \cite{Kirchner16,Pfirrmann16}) or by modifying PMMA  through a process such as \emph{sequential infiltration synthesis} \cite{Tseng11}. 

Overall, TASTE provides an alternative approach to fabricating blazed gratings that perform with high diffraction efficiency at soft x-ray wavelengths. 
While \ce{KOH} etching may be capable of producing gratings that perform with higher diffraction efficiency owing to the smooth facets and well-defined structures generated by the process [\emph{cf.\@} \cref{sec:crystal_etching}], TASTE has the key advantage for variable-line-space gratings that the precision of $d$ is not limited by the crystal structure of the substrate [\emph{cf.\@} \cref{fig:crystal_structure}]. 
This is expected to be beneficial for the radially-ruled gratings that the \emph{XGS} reflection grating design calls for \cite{McEntaffer19}, but due to the relatively low throughput of TASTE, like other EBL-based processes, the implementation of a grating replication process is required [\emph{cf.\@} \cref{sec:SCIL_ch5}]. 
Moreover, the \emph{Extreme-Ultraviolet Stellar Characterization for Atmospheric Physics and Evolution (ESCAPE)} mission concept, with a goal of characterizing high-energy radiation in habitable zones surrounding M-dwarfs and their impact on the atmospheres of exoplanets, baselines a spectrometer design that incorporates two large-area blazed gratings featuring curved groove layouts that play a similar role for extreme ultraviolet (EUV) radiation in a \emph{Hettrick-Bowyer-I} telescope \cite{France19,France20,Hettrick84}. 
With \emph{ESCAPE} baselining $\sim \SI{60}{\percent}$ single-order diffraction efficiency in the EUV, the beamline results presented in \cref{sec:taste_eff_results} give an indication that TASTE is capable of meeting this goal for spectral sensitivity. 

Further work involving beamline testing for $\mathscr{R}$ is required to determine whether TASTE is capable of producing gratings that meet performance requirements for instruments such as \emph{ESCAPE} and the \emph{XGS} for \emph{Lynx} \cite{France19,McEntaffer19}. 
Smaller-scale development for efficient, high-resolution, x-ray reflection gratings is taking place for the \emph{The Off-plane Grating Rocket Experiment (OGRE)} \cite{DeRoo13,McEntaffer14,Donovan18b,Tutt18,Donovan19a,Donovan19b}, which calls for $\delta \lessapprox 30^{\circ}$ blazed gratings with $d \lessapprox \SI{160}{\nm}$ radial-groove layouts that match a \SI{3.5}{\metre} focal length. 
Along with optical-system testing for $\mathscr{R}$, grating development for \emph{ESCAPE} and \emph{OGRE} warrants further diffraction-efficiency testing to determine which blazed-grating manufacture produces curved or fanned grooves that are efficient while also offering high $\mathscr{R}$. 
A new ALS general user proposal to carry out diffraction-efficiency experimentation for \emph{ESCAPE} and \emph{OGRE} grating prototypes has recently been submitted for beamline testing in late 2021 and 2022. 
These prototypes include four grating variants, all with the same groove layout: one fabricated via EBL coupled with \ce{KOH} etching \cite{Miles18}, one via TASTE \cite{McCoy20,McCurdy20}, one via EBL coupled with directional ion milling \cite{Miles2019_ion_mill} and one analogous, laminar grating for comparison. 
Together with beamline results for $\mathscr{R}$, these results will inform which fabrication technique is best suited for achieving high efficiency and $\mathscr{R}$ simultaneously. 

%which seeks to observe emission lines that emanate from coronal activity in the Capella star system

%Further work involving beamline testing for $\mathscr{R}$, however, is required to determine if TASTE can produce gratings that perform with $\mathscr{R} \gtrapprox 5000$ in the converging beam of a Wolter-I telescope. 

\section{SCIL for Grating Replication}\label{sec:SCIL_ch5}
%%%%%%%%%%%%%%%%%%%%%%%%%%%%%%%%%%%%%%%%%--------------------------------------------------
SCIL is a variant of \emph{nanoimprint lithography (NIL)} designed for high-throughput patterning of nanoscale structures over large areas \cite{Verschuuren10,Verschuuren17}. 
%As described in \cref{sec:grating_fab}, 
The process utilizes a flexible stamp that is molded from a rigid master template to imprint features in resist with the aid of specialized pneumatic tooling [\emph{cf.\@} \cref{sec:grating_fab}]. 
In contrast to rigid-stamp UV-NIL, which has been pursued previously for x-ray reflection grating replication \cite[\emph{cf.\@} \cref{sec:nanoimprint}]{Chang03,Chang04,McEntaffer13,DeRoo16,Miles18}, SCIL enables conformal imprinting over wafers up to \SI{200}{\mm} in diameter without the need for high applied pressure. 
This in turn avoids damage to an expensive master template while reducing pattern defects and trapped air pockets. % and, 
Moreover, packaged equipment developed by \textsc{Philips SCIL Nanoimprint Solutions} \cite{philis_scil}, known commercially as \textsc{AutoSCIL}, serves to automate the imprinting process for high-volume production, which is most compatible with \textsc{NanoGlass}, an inorganic resist that cures through a thermodynamically-driven, silica sol-gel process \cite{Verschuuren10,Verschuuren17,Verschuuren19}. 
Using \textsc{AutoSCIL}, a single stamp is capable of producing $\gtrapprox \num{700}$ imprints in \textsc{NanoGlass} without pattern degradation, at a rate of \num{60} wafers per hour \cite{Verschuuren17,Verschuuren18}. 

In collaboration with \textsc{Philips SCIL Nanoimprint Solutions}, the \textsc{AutoSCIL} technique was first applied to x-ray reflection grating technology for the \emph{Water Recovery X-ray Rocket (WRXR)}, which was a Penn State sounding-rocket payload that utilized \num{26} replicas of a \ce{KOH}-etched silicon master for a diffuse-object, soft x-ray spectrometer \cite{Miles17,Miles18b,Miles19,Verschuuren18,Tutt19}. 
This production run served as a trial, indicating that the process is well suited for producing many replicas of a master grating for the purpose of populating modular arrays in a Wolter-I spectrometer, which is of particular importance for \emph{XGS} as it calls for thousands of identical gratings \cite{McEntaffer19}, in addition to upcoming rocket payloads that each call for hundreds of replicas: \emph{The Rockets for Extended-source X-ray Spectroscopy (tREXS)} \cite{Miles19b,Tutt19b} and \emph{OGRE} \cite{Tutt18,Donovan19a,Donovan19b}. 
As in virtually any NIL process, however, imprinted surface-relief features are prone to topographic distortion induced by resist shrinkage \cite{Horiba_2012}. 
This phenomenon occurs in \textsc{NanoGlass} as the silica precursors tetramethylorthosilicate (TMOS) and methyltrimethoxysilane (MTMS) contained in the resist [\emph{cf.\@} \cref{fig:precursors}] react to form a silica-like network with a cross-link density that depends on the post-imprint cure temperature, $T_{\text{cure}}$ \cite{Verschuuren10}. 
In \cref{ch:grating_replication}, it is demonstrated through diffraction-efficiency testing at beamline 6.3.2 of the ALS \cite{ALS_632,McCoy20b} that $T_{\text{cure}} \approx \SI{90}{\celsius}$ gives rise to a $\sim \SI{10}{\percent}$ volumetric shrinkage in the resist, which causes a $\sim 2^{\circ}$ reduction in $\delta$ relative to a $\sim 30^{\circ}$ facet angle defined by \ce{KOH} etching in $\langle 311 \rangle$-oriented silicon [\emph{cf.\@} \cref{fig:KOH_geo_311}]. 
With support from AFM measurements, this was carried out by comparing the blaze response of a test imprint provided by \textsc{Philips SCIL Nanoimprint Solutions} [\emph{cf.\@} \cref{fig:SCIL_SEMs}] and the corresponding \ce{KOH}-etched silicon master [\emph{cf.\@} \cref{fig:master_grating_SEM}]. %, which was fabricated by staff at the Penn State Materials Research Institute [\emph{cf.\@} \cref{fig:master_grating_SEM}]. 

The result just described shows that the impact of resist shrinkage on blaze angle is non-negligible and hence should be compensated for in the fabrication of a master grating \cite{McCoy20b}. 
Further studies to examine how the reduced blaze angle, $\delta'$, evolves with $T_{\text{cure}}$ should aid in this process, where fabricating a master grating with a precise value for $\delta$ may prove to be difficult. 
This is expected to be useful for achieving a blaze angle that is intermediate between what is is produced from \ce{KOH} etching in silicon with standard wafer orientations [\emph{cf.\@} \cref{tab:off_axis_si}]. 
Although the \textsc{AutoSCIL} production platform provides an avenue for high-volume production of grating imprints, PVD techniques such as EBPVD and \emph{plasma sputter deposition} are limited in throughput, and moreover, the impact of ion bombardment (from the latter process) on the sol-gel network has not yet been investigated. 
This motivates the pursuit of deposition strategies that are both capable of high throughput and compatible with \textsc{NanoGlass} resist. 

In addition to studying how $\delta'$ depends on $T_{\text{cure}}$ for a given groove geometry imprinted in resist and the effects of PVD coatings on groove shape, film-stress studies are warranted to determine how depositions of \textsc{NanoGlass} resist and metallic, reflective coatings contribute to substrate deformation \cite{Scheuerman69}. 
Similar to studies that examine how film stress affects the performance of x-ray telescope optics in terms of angular resolution \cite{Chan13,Probst18,Molnar-Fenton19}, compensating for this effect is motivated by the drive to preserve the flat figure of each grating in a spectrometer so that $\mathscr{R}$ is not degraded by unintentional substrate curvature. 
With the grooves of grating replicas being carried in \textsc{NanoGlass}, the nanoscale porosity of the resist network caused by the MTMS precursor is prone to the entrapment of water vapor in a manner similar to other sol-gel systems \cite{Hench90,Brinker90,Vishnevskiy19}. 
While it has been found experimentally that water present in the resist network can lead to minor imprinting issues that can be overcome \cite{Verschuuren10}, the effect of water vapor on aged imprints coated for reflectivity should be further investigated. 
Alternatively, the sol-gel network can be fully densified with $T_{\text{cure}} \gtrapprox \SI{850}{\celsius}$ to eliminate porosity but this comes with increased film stress as well as a smaller value for $\delta'$. % and potential groove-shape deformation. 

%stress studies and effect of water vapor on sol-gel network
%[To be expanded]

The nanofabrication processes of TASTE and SCIL, together, offer an avenue for manufacturing next-generation x-ray reflection gratings. 
With established process development for TASTE and equipment for SCIL stamp construction as well as low-volume, pneumatic imprinting (built by \textsc{S{\"U}SS MicroTec} \cite{suss_microtec}) recently installed at the Penn State Nanofabrication Laboratory \cite{PSU_MRI_nanofab}, future work will include studying the compatibility of SCIL with a master grating fabricated by TASTE.  %(using a mask aligner add-on built by \textsc{S{\"U}SS MicroTec})
Because polydimethylsiloxane (PDMS) [\emph{cf.\@} \cref{fig:PDMS}] and PMMA do not bond well without the use of an adhesive layer \cite{Tan10}, it is hypothesized that an anti-stiction treatment is not needed to construct an X-PDMS SCIL stamp from TASTE-processed resist; this is supported by a study that shows how PDMS and H-PDMS stamps used in \emph{soft UV-NIL} processes can be molded directly from untreated PMMA patterned by EBL \cite{Huelsen14}. As oxidized silicon is exposed between groove facets in a TASTE grating, however, it is likely that a small level of stiction will occur and hence stamp separation should be performed carefully.
The processes of TASTE and SCIL are nevertheless expected to be compatible for the production of grating replicas. 
This will be pursued experimentally using a duplicated, uncoated version of the TASTE prototype described in \cref{sec:TASTE_prototype}. 

\section{Outlook for Future Studies}\label{sec:summary_ch5}
%%%%%%%%%%%%%%%%%%%%%%%%%%%%%%%%%%%%%%%%%--------------------------------------------------
With its ability to generate blazed groove facets in resist over a custom layout defined by EBL, TASTE is a promising technique for the manufacture of state-of-the-art x-ray reflection gratings \cite[\emph{cf.\@} \cref{sec:TASTE_ch5}]{McCoy18,McCoy20}. 
The process does, however, face limitations in its ability to produce sharply-defined groove facets as the groove periodicity, $d$, becomes significantly smaller than \SI{400}{\nm} such that the width of each GEBL staircase step is comparable to the diameter of the focused electron beam \cite{McCurdy20}. 
While pattern resolution has potential for improvement through the implementation of a cold-development process \cite{Ocola06}, a fundamental size-scale limitation likely exists for GEBL \cite{Fallica17}. 
This motivates the use of \emph{EUV lithography} \cite{Wu07,Tallents10} or \emph{EUV interference lithography (EUV-IL)} \cite{Auzelyte09} at $\lambda \approx \SI{13.5}{\nm}$, which are processes capable of patterning on the sub-\SI{10}{\nm} level. 
Of particular interest is grayscale EUV-IL, which has been demonstrated to be capable of generating a $d = \SI{100}{\nm}$ tri-level staircase pattern in PMMA with the use of a transmission-grating EUV exposure mask that provides a two-fold reduction in pattern periodicity and a mask-shifting technique that enables dose-modulated exposure \cite{Fallica17,Buitrago16}. 
Such a process can, in principle, be coupled with thermal reflow to realize an EUV-IL variant of the TASTE process to produce a surface relief for a blazed x-ray reflection \cite{Fallica17,Schleunitz14}. 
Further investigation, however, is required to determine if this fabrication approach can produce a large-area groove layout with a high-fidelity radial profile that also enables high diffraction efficiency. 
This would necessitate the manufacture of a transmission-grating mask that is designed appropriately so as to produce a radial-groove diffraction pattern and additionally, the use of projection lithography equipment with a high-brightness source of EUV radiation (\emph{e.g.}, beamline 12.0.1 of the ALS \cite{ALS_EUVlitho}). 

This dissertation has considered only the fabrication of planar gratings; \emph{concave gratings} \cite{Loewen97}, however, are also of interest for their ability to focus and disperse radiation simultaneously. 
The fabrication of gratings on curved surfaces has been pursued through a variety of techniques throughout the years, including mechanical ruling \cite{Harada80,Kita83,Kita92}, holographic recording \cite{Sokolova04}, directional ion etching \cite{Liu15,Guo18}, soft lithography \cite{Xia98,Paul03,Choi04}, projection lithography \cite{Klosner02}, as well as hybrid approaches that combine aspects of soft and projection lithography \cite{Kim09,Park12}. 
Additionally, EBL is capable of patterning on curved substrates \cite{Wilson03,PSU_MRI_EBL,Lee18,Arat19} and this is of special interest in x-ray spectroscopy primarily for the following reason. 
That is, if a blazed x-ray grating topography can be patterned on the surface of a hyperbolic mirror appropriately, a two-element Wolter-I grating spectrometer could be realized \cite{DeRoo19,DeRoo19b}. 
Such an optical system would reduce the number of reflections required to produce x-ray spectra while reducing instrument mass and eliminating the need for grating-array alignment. 
However, strategies for the mass production of hyperbolic gratings should be considered. 
This may involve using a SCIL stamp to pattern on curved surfaces in a manner similar to soft-lithographic approaches \cite{Paul03,Choi04} but further research and experimentation is needed to make this assessment. 

Beyond telescopic applications of concave gratings, the nanofabrication technology discussed in this thesis is also of interest for laboratory x-ray astrophysics experiments \cite{Beiersdorfer03,Hell20} that utilize an \emph{electron-beam ion trap (EBIT)} \cite{Levine88a,Levine88b} to produce spectra from quasi-stationary, highly-charged ions. 
Such an x-ray source effectively provides diverging rays from a central electron beam where a particular species of ion is trapped; this radiation must be focused and dispersed to produce spectra with appreciable $\mathscr{R}$ and high diffraction efficiency. % is desirable for spectral sensitivity. 
While commercially-available, mechanically-ruled concave gratings with in-plane, variable-line-space layouts have been implemented for EUV and soft x-ray spectroscopy on EBIT systems previously \cite{Beiersdorfer99,Utter99,Beiersdorfer04,Ohashi11,Shi14}, these instruments achieve at most $\mathscr{R} \approx 1200$ for long-$\lambda$ soft x-rays. 
This performance can, in principle, be improved by constructing a new EBIT grating spectrometer that is centered around a custom, blazed grating patterned on an ellipsoidal mirror segment that brings diverging rays to a focus while also producing an off-plane diffraction pattern that enables high-$\mathscr{R}$ spectra to be extracted through the use of high diffracted orders [\emph{cf.\@} \cref{sec:grating_tech_intro}]. 
Owing to its ability to pattern a blazed grating topography with the precision of EBL, TASTE is hypothesized to be capable of fabricating such a custom, concave grating provided that an appropriate curved substrate that focuses radiation effectively can be obtained. 
The development of a next-generation grating spectrometer for EBIT spectroscopy and its installment at an EBIT laboratory, such as those at Lawrence-Livermore National Laboratory \cite{llnl_ebit} and the Harvard-Smithsonian Center for Astrophysics \cite{cfa_ebit}, is crucial for furthering the fields of astrophysical x-ray spectroscopy and theoretical atomic physics by enabling wavelength centroids, transition rates and cross-sections of faint, poorly-studied, soft x-ray spectral lines to be measured in the laboratory so that they can complement astrophysical observations \cite{Smith19_labastro}. 
At the time of this writing, preliminary research is being undertaken to move forward with this project. 
%rather than relying on photon-starved observations of astrophysical plasmas


%EBIT \cite{Levine88a,Levine88b,Beiersdorfer03,Smith19_labastro}

%[To be expanded]

